\documentclass[11pt]{article}

\usepackage[T1]{fontenc}
\usepackage{amsmath,amssymb,amsthm}
\usepackage[margin=1in]{geometry}
\usepackage[hidelinks,pdfauthor={Xiaojian Zeng},pdftitle={Polynomial Solvability of a Two-Machine Unit-Time Resource-Constrained Job Shop with at Most Three Operations per Job},pdfsubject={Scheduling complexity and dynamic programming},pdfkeywords={job shop scheduling, resource constraints, unit processing times, dynamic programming}]{hyperref}

\newtheorem{theorem}{Theorem}
\newtheorem{lemma}[theorem]{Lemma}
\newtheorem{remark}[theorem]{Remark}

\newcommand{\Cmax}{C_{\max}}
\newcommand{\N}{\mathbb{N}}

\title{Polynomial Solvability of a Two-Machine Unit-Time Resource-Constrained Job Shop with at Most Three Operations per Job}
\author{Xiaojian Zeng\\
PBC School of Finance, Tsinghua University, 100190\\
\texttt{xiaojianzeng@foxmail.com}}
\date{}

\begin{document}
\maketitle

\begin{abstract}
We consider two dedicated machines, nonpreemptive unit-time operations, at
most three operations per job in an alternating prescribed route, and one
unit-capacity auxiliary resource.  Every operation of a job has the same
binary resource demand.  We prove that minimum makespan is polynomial-time
solvable.  After an integral-slot normalization, a job is characterized by
one of twelve route--demand types and by the number of operations it has
completed.  A shortest-path dynamic program over the counts in these prefix
classes exactly represents all feasible schedules.  Its state space has at
most $(n+3)^{24}$ states, and an optimal labeled schedule can be recovered.
This settles the at-most-three-operation boundary under the stated model; it
does not settle the corresponding problem with unbounded route length, and
the exponent is not intended as a practical running-time bound.
\end{abstract}

\section{Context and contribution}

Tellache and Azerine study a job shop with a conflict graph and its
equivalent formulation using unit-capacity auxiliary resources whose demands
are constant over each job \cite{TellacheAzerine2026}.  In the two-machine,
unit-processing-time setting they prove polynomial solvability when every job
has at most two operations, and they leave the version with unrestricted
numbers of operations unresolved.  The result here treats the intervening
boundary of at most three operations per job.

It is important not to infer hardness at this boundary from the
resource-free problem with arbitrary processing times.  The classical
two-machine, at-most-three-operation NP-completeness construction of Lenstra,
Rinnooy Kan, and Brucker uses processing times derived from a \textsc{Knapsack}
instance \cite{LenstraRinnooyKanBrucker1977}.  It therefore does not apply to
the all-unit-time restriction.  Conversely, without the auxiliary resource,
the two-machine unit-operation-time job shop is solvable in linear time
\cite{HefetzAdiri1982}.  The demand-one jobs in the present problem form an
additional clique of pairwise nonoverlap constraints, so that resource-free
algorithm does not itself resolve the model considered here.

The proof below is self-contained.  Its contribution relative to
\cite{TellacheAzerine2026} is a finite-type, aggregate-prefix dynamic program
that advances the printed bound from two operations to three.  A bounded
search of exact notation and equivalent conflict-graph formulations did not
locate the same at-most-three-operation statement.  This observation is only
a description of the sources inspected, not a claim of priority.

\section{Model and main result}

There are two dedicated machines, denoted $A$ and $B$, and a finite set of
jobs $J=\{J_1,\ldots,J_n\}$.  Job $J_j$ consists of a chain
\[
 O_{1j}\prec O_{2j}\prec\cdots\prec O_{m_jj},
 \qquad 1\leq m_j\leq 3.
\]
Each operation is nonpreemptive, has processing time one, and is assigned to
one of the two machines.  Consecutive operations of a job use different
machines; a machine may be revisited.  Each job has a demand bit
$r_j\in\{0,1\}$, and every one of its operations has that same demand for a
single renewable auxiliary resource of capacity one.

A schedule assigns a real start time $s_{ij}\geq0$ to every operation.
Intervals are taken half-open.  The operations of a job respect their chain,
operations assigned to the same machine do not overlap, and no two operations
with demand one overlap.  There are no release dates, setup times, or further
resources.  The objective is to minimize
\[
 \Cmax=\max_{i,j}(s_{ij}+1).
\]

\begin{theorem}\label{thm:main}
For every instance of the model above, an optimal makespan and an optimal
labeled schedule can be computed in polynomial time.  In particular, a
shortest-path dynamic program with at most $(n+3)^{24}$ states solves
\[
 J2\mid \mathrm{res}^{t}111,\ p_{ij}=1,\ m_j\leq3\mid \Cmax.
\]
\end{theorem}

The superscript $t$ records that the resource demand is constant over all
operations of a job.  We first justify the discrete time model used by the
dynamic program.

\section{Integral-slot normalization}

\begin{lemma}[Integral slots]\label{lem:integral}
Every feasible schedule has a feasible schedule with integer start times and
no larger makespan.  Moreover, an optimal integer schedule can be chosen so
that every unit slot before its completion contains at least one operation.
\end{lemma}

\begin{proof}
Start from an arbitrary feasible schedule and make one vertex for every
operation.  Include all prescribed job-precedence arcs.  For every pair of
operations that cannot overlap---because they use the same machine, both
have demand one, or belong to the same job---orient the pair in the temporal
order realized by the schedule.  Thus every resulting arc $u\to v$ satisfies
\[
 s_v\geq s_u+1.
\]
The digraph is acyclic: a directed cycle of $q>0$ arcs would imply
$s_u\geq s_u+q$ for one of its vertices.

For a vertex $v$, let $h(v)$ be the maximum number of arcs on a directed path
ending at $v$; a source has value zero.  If $u\to v$, then
$h(v)\geq h(u)+1$.  Hence assigning the new start time $h(v)$ preserves every
precedence and nonoverlap constraint.  These start times are nonnegative
integers.  A path of $h(v)$ arcs ending at $v$ already implies
$s_v\geq h(v)$ in the original schedule, and consequently
\[
 h(v)+1\leq s_v+1\leq \Cmax.
\]
The new makespan is therefore no larger.

Finally, if an integer schedule has a globally empty slot, shift every
operation strictly after that slot one unit to the left.  Unit processing
times and integral starts imply that no operation crosses the empty slot, so
the shift preserves all constraints.  Repeating this deletion removes every
empty slot without increasing the makespan.
\end{proof}

\section{Aggregate-prefix dynamic program}

Because consecutive machines alternate, the possible routes are
\[
 A,\quad B,\quad AB,\quad BA,\quad ABA,\quad BAB.
\]
Crossing these six routes with $r\in\{0,1\}$ gives a fixed set $T$ of twelve
types.  For $t\in T$, let $\ell_t\in\{1,2,3\}$ be the route length and let
$n_t$ be the number of jobs of that type.

A state is the vector
\[
 x=(x_{t,k}:t\in T,\ 0\leq k\leq\ell_t),
\]
where $x_{t,k}$ is the number of type-$t$ jobs that have completed exactly
$k$ operations.  Thus $\sum_{k=0}^{\ell_t}x_{t,k}=n_t$.  Initially
$x_{t,0}=n_t$ for every $t$; at the terminal state
$x_{t,\ell_t}=n_t$.

A prefix class $(t,k)$ is ready when $k<\ell_t$ and $x_{t,k}>0$.  Its next
operation has the machine specified by position $k+1$ of the route and the
demand bit of $t$.  A slot action selects either one ready class or two
distinct ready classes.  In the latter case the two next operations must use
different machines and the sum of their demand bits must be at most one.  For
each selected class, a transition decreases $x_{t,k}$ by one and increases
$x_{t,k+1}$ by one.  Every transition has cost one.  The state graph is
acyclic because the total number of completed operations increases at every
transition.

\begin{lemma}[Exact quotient]\label{lem:quotient}
There is a path of $q$ transitions from the initial state if and only if
there is a realizable sequence of $q$ nonempty integer slots reaching the
same aggregate state.  Both implications preserve the length.
\end{lemma}

\begin{proof}
Consider first a transition.  Choose an arbitrary labeled job from every
selected prefix class and run its next operation in the new slot.  Distinct
prefix classes contain disjoint sets of jobs, the chosen machines are
different, and at most one chosen operation has demand one.  All job,
machine, and resource constraints therefore hold.  Induction realizes every
DP path as a labeled schedule of the same length.

Conversely, consider a feasible nonempty integer slot and the aggregate state
immediately before it.  At most one operation runs on each of the two
machines.  Every running operation is the ready next operation of its job,
the operations use distinct machines, and at most one has demand one.  The
slot therefore selects exactly an allowed set of one or two prefix classes
and induces the corresponding transition.  Reading the slots in order gives
a DP path.  The aggregate vector loses only the identities of interchangeable
jobs: type and completed prefix determine all information relevant to their
next operation.
\end{proof}

\begin{proof}
By Lemma~\ref{lem:integral}, some optimum is an integer schedule with no
empty slot.  Lemma~\ref{lem:quotient} maps this schedule to a DP path of the
same length.  Its reverse direction realizes every DP path as a feasible
schedule.  A shortest initial-to-terminal path therefore has length equal to
the optimal makespan.

For a fixed type $t$, the number of nonnegative prefix distributions is
\[
 \binom{n_t+\ell_t}{\ell_t}
 \leq (n+3)^{\ell_t}.
\]
The sum of the route lengths over all twelve types is
\[
 \sum_{t\in T}\ell_t
 =2(1+1+2+2+3+3)=24.
\]
Hence the total number of states is at most $(n+3)^{24}$.  There are at most
24 ready prefix classes, so a state has fewer than
$24+\binom{24}{2}$ candidate actions.  State generation and shortest-path
optimization therefore take polynomial time in the explicit input size.
To recover labels, maintain a list of jobs in each prefix class and record the
chosen job whenever a transition moves one list element to the next class.
\end{proof}

\section{The zero-demand inclusion and the boundary of the result}

The printed resource convention in \cite{TellacheAzerine2026} permits zero
demands and explicitly uses both zero- and one-demand jobs.  Given an instance
$I$ of the resource-free two-machine unit-time job shop with the same route
bound, retain all jobs, routes, and processing times and set $r_j=0$ for every
job.  A schedule is feasible before this map if and only if it is feasible
after it: the auxiliary-resource load is identically zero.  The map preserves
all start times, makespans, and optima.

If an additional convention requires a nonempty instance to contain at least
one demand-one job, no dummy job is needed.  Select one existing job, give it
demand one, and leave every other demand at zero.  That job's own operations
are already sequential, and no other job uses the resource, so the added
resource condition remains vacuous.  This constant-overhead map again
preserves every schedule and its makespan.  By contrast, adding a new
one-operation dummy can consume a machine slot and need not preserve the
optimum, even though the dummy has bounded size.

This inclusion does not furnish hardness.  In the reduction of
\cite{LenstraRinnooyKanBrucker1977}, the relevant two-machine jobs have
variable processing times such as values taken from the \textsc{Knapsack}
instance and the quantities $b$ and $A-b$.  Thus the reduction neither
applies to $p_{ij}=1$ nor proves strong NP-hardness of the unit-time
restriction.  The polynomial result of \cite{HefetzAdiri1982} is consistent
with this distinction.

Theorem~\ref{thm:main} relies essentially on a constant set of route--prefix
classes.  When route length is unbounded, the number of such classes is not
fixed and the exponent in the state-count argument is no longer constant.
The theorem consequently makes no complexity claim for the unrestricted
operation-count problem posed in \cite{TellacheAzerine2026}.  Even in the
three-operation case, the exponent 24 is deliberately crude; the result is a
complexity classification, not a claim of practical efficiency.

\section{Exact check and assistance disclosure}

As a diagnostic, two separately coded exact recursions were compared.  One
uses only the aggregate prefix counts above; the other retains every job
label.  The program exhausts all
\[
 \binom{6+12}{12}=18{,}564
\]
weak twelve-type count vectors containing at most six jobs, including the
empty vector.  The two recursions returned equal optimal makespans in every
case.  The calculation uses exact integer states, no random choices, and no
floating-point comparisons.  From the research-artifact directory, it can be
reproduced with
\begin{center}
\texttt{python3 evidence/aggregate\_dp\_check.py -{}-max-jobs 6}.
\end{center}

Independent exact implementations are useful for reducing implementation
risk in finite verification; for a recent example using integer predicates
and two exact maximum-clique implementations, see
\cite{ZengEtAl2026Certificates}.  In the present paper the finite comparison
has the more limited role of a sanity check.  It is not a premise of
Theorem~\ref{thm:main}; the integral-slot lemma and the two directions of
Lemma~\ref{lem:quotient} carry the proof for arbitrary $n$.

OpenAI Codex assisted with the bounded literature search, implementation of
the two exact recursions, adversarial proof checking, and preparation of the
manuscript.  Language-model output is not treated as mathematical evidence.
No proof assistant was used.

\bibliographystyle{abbrv}
\bibliography{references}

\end{document}
